\subsection{Logische Verknüpfungen}
\label{subsec:fundamentals_logical_connectives}

Durch die Verknüpfung zweier Aussagen $A$ und $B$ lassen sich neue Aussagen bilden:

\begin{DefinitionBox}{UND-Verknüpfung $\land$}{fundamentals_and_connective}
  Die \emph{UND-Verknüpfung}, fachsprachlich \emph{Konjunktion}
  \EnglishTerm{conjunction},
  \[
    A\land B,
  \]
  ist nur dann wahr, wenn sowohl $A$ als auch $B$ wahr sind.
\end{DefinitionBox}

\begin{DefinitionBox}{ODER-Verknüpfung}{fundamentals_or_connective}
  Die \emph{ODER-Verknüpfung}, fachsprachlich \emph{Disjunktion}
  \EnglishTerm{disjunction},
  \[
    A\lor B,
  \]
  ist wahr, wenn mindestens eine der beiden Aussagen $A$ und $B$ wahr ist.
\end{DefinitionBox}

\begin{samepage}
Eine Wahrheitstabelle erfasst alle möglichen Wahrheitswerte von $A$ und $B$;
durch das Ausfüllen erkennen wir systematisch, wann die UND- und die
ODER-Verknüpfung wahr oder falsch sind.

\begin{center}
  \renewcommand{\arraystretch}{1.55}
  \setlength{\tabcolsep}{12pt}
  \begin{tabular}{cccc}
    \toprule
    $A$ & $B$ & $A\land B$ & $A\lor B$\\
    \midrule
    T & T & \Cloze{1.8cm}{T} & \Cloze{1.8cm}{T}\\
    T & F & \Cloze{1.8cm}{F} & \Cloze{1.8cm}{T}\\
    F & T & \Cloze{1.8cm}{F} & \Cloze{1.8cm}{T}\\
    F & F & \Cloze{1.8cm}{F} & \Cloze{1.8cm}{F}\\
    \bottomrule
  \end{tabular}
\end{center}
\end{samepage}

\begin{RemarkBox}[remark title=Inklusives und exklusives ODER]
  Das mathematische ODER ist grundsätzlich inklusiv, d.\,h. $A\lor B$ ist auch
  dann wahr, wenn sowohl $A$ als auch $B$ wahr sind. Im Gegensatz dazu gibt es
  in der Informatik auch das exklusive ODER, das nur dann wahr ist, wenn genau
  eine (aber nicht beide) der Aussagen wahr sind.
\end{RemarkBox}
